\documentclass[11pt,a4paper]{article}
\usepackage[margin=1.8cm,top=2.5cm,bottom=2.2cm,headheight=24pt]{geometry}
\usepackage{fontspec}
\setmainfont{RakutenSerifApp}[
  Path = /home/sudheesh/.fonts/Notesnook/,
  Extension = .ttf,
  UprightFont = *-Regular,
  BoldFont = *-Bold,
  ItalicFont = *-Italic,
  BoldItalicFont = *-BoldItalic
]
\usepackage{amsmath,amssymb,mathtools}
\usepackage{booktabs,array,tabularx,multirow}
\usepackage{graphicx}
\usepackage{xcolor}
\usepackage{titlesec}
\usepackage{fancyhdr}
\usepackage{enumitem}
\usepackage{caption}
\usepackage{tcolorbox}
\tcbuselibrary{skins,breakable}
\usepackage[hidelinks,pdfusetitle]{hyperref}

% Define custom column for tabularx
\newcolumntype{Y}{>{\centering\arraybackslash}X}

% Define color palette
\definecolor{navyprimary}{RGB}{15, 35, 65}
\definecolor{accentblue}{RGB}{2, 132, 199}
\definecolor{tealaccent}{RGB}{13, 148, 136}
\definecolor{cardbg}{RGB}{248, 250, 252}
\definecolor{cardborder}{RGB}{203, 213, 225}
\definecolor{calloutbg}{RGB}{241, 245, 249}
\definecolor{calloutframe}{RGB}{100, 116, 139}
\definecolor{formulabg}{RGB}{238, 242, 255}
\definecolor{formulaframe}{RGB}{99, 102, 241}
\definecolor{examplebg}{RGB}{250, 250, 250}
\definecolor{exampleframe}{RGB}{56, 189, 248}
\definecolor{highlightgold}{RGB}{180, 83, 9}
\definecolor{alertbg}{RGB}{254, 242, 242}
\definecolor{alertframe}{RGB}{239, 68, 68}

% Caption styling
\captionsetup{font=small,labelfont={bf,color=navyprimary},skip=5pt}

% Typography & Header/Footer
\pagestyle{fancy}
\fancyhf{}
\renewcommand{\headrulewidth}{0.6pt}
\renewcommand{\headrule}{\hbox to\headwidth{\color{navyprimary}\leaders\hrule height \headrulewidth\hfill}}
\fancyhead[L]{\small\bfseries\color{navyprimary}Power Electronics (Module 2) --- Master Study Guide}
\fancyhead[R]{\small\color{accentblue}CET S5 EEE}
\fancyfoot[L]{\small\color{gray}College of Engineering Trivandrum}
\fancyfoot[R]{\small\bfseries\color{navyprimary}Page \thepage}

% Section Titles Styling
\titleformat{\section}{\Large\bfseries\color{navyprimary}}{}{0em}{}[\vspace{1mm}{\color{accentblue}\hrule height 1.2pt}]
\titleformat{\subsection}{\large\bfseries\color{accentblue}}{}{0em}{}
\titleformat{\subsubsection}{\normalsize\bfseries\color{tealaccent}}{}{0em}{}

% Tcolorbox custom environments (Border-free modern cards with solid background colors)
\newtcolorbox{conceptbox}[1][]{
  enhanced,
  colback=calloutbg,
  colframe=calloutbg,
  colbacktitle=calloutbg,
  coltitle=navyprimary,
  fonttitle=\bfseries,
  title={#1},
  titlerule=0pt,
  boxrule=0pt,
  frame hidden,
  arc=2.5mm,
  left=4mm, right=4mm, top=3mm, bottom=3mm,
  breakable
}

\newtcolorbox{formulabox}[1][]{
  enhanced,
  colback=formulabg,
  colframe=formulabg,
  colbacktitle=formulabg,
  coltitle=navyprimary,
  fonttitle=\bfseries,
  title={#1},
  titlerule=0pt,
  boxrule=0pt,
  frame hidden,
  arc=2.5mm,
  left=4mm, right=4mm, top=3mm, bottom=3mm,
  breakable
}

\newtcolorbox{examplebox}[1][]{
  enhanced,
  colback=examplebg,
  colframe=examplebg,
  colbacktitle=examplebg,
  coltitle=navyprimary,
  fonttitle=\bfseries,
  title={#1},
  titlerule=0pt,
  boxrule=0pt,
  frame hidden,
  arc=2.5mm,
  left=4mm, right=4mm, top=3mm, bottom=3mm,
  middle=2.5mm,
  segmentation style={solid, draw=cardborder, line width=0.6pt},
  breakable
}

\newtcolorbox{alertbox}[1][]{
  enhanced,
  colback=alertbg,
  colframe=alertbg,
  colbacktitle=alertbg,
  coltitle=alertframe,
  fonttitle=\bfseries,
  title={#1},
  titlerule=0pt,
  boxrule=0pt,
  frame hidden,
  arc=2.5mm,
  left=4mm, right=4mm, top=3mm, bottom=3mm,
  breakable
}

% Remove black borders from \boxed formulas while keeping their prominent display
\renewcommand{\boxed}[1]{#1}


% Float placement parameters to prevent large blank spaces
\renewcommand{\topfraction}{0.90}
\renewcommand{\bottomfraction}{0.85}
\renewcommand{\textfraction}{0.10}
\renewcommand{\floatpagefraction}{0.75}
\setcounter{topnumber}{4}
\setcounter{bottomnumber}{4}
\setcounter{totalnumber}{8}
\raggedbottom

\begin{document}

% Title Header
\begin{center}
  {\huge\bfseries\color{navyprimary} POWER ELECTRONICS (PE) --- MODULE 2}\\[1.5ex]
  {\Large\bfseries\color{accentblue} The Definitive Master Study Guide, Waveform Atlas \& Derivation Handbook}\\[1.5ex]
  {\normalsize\bfseries Course: B.Tech S5 Electrical \& Electronics Engineering (EEE) \quad|\quad College of Engineering Trivandrum (CET)}\\[0.8ex]
  {\small Integrated References: Dr. P. S. Bimbhra, Ned Mohan, Daniel W. Hart, L. Umanand \& Official Course Lecture Waveforms}
\end{center}
\vspace{-1em}
\noindent\textcolor{accentblue}{\rule{\textwidth}{1.5pt}}
\vspace{0.5em}

\tableofcontents
\newpage

%%%%%%%%%%%%%%%%%%%%%%%%%%%%%%%%%%%%%%%%%%%%%%%%%%%%%%%%%%%%%%%%%%%%%%%%%%%%%%%
\section{Module 2 Syllabus, Architecture \& Learning Roadmap}
%%%%%%%%%%%%%%%%%%%%%%%%%%%%%%%%%%%%%%%%%%%%%%%%%%%%%%%%%%%%%%%%%%%%%%%%%%%%%%%

According to the official APJ Abdul Kalam Technological University / CET syllabus, Module 2 carries substantial weightage in both Continuous Internal Evaluation (CIE) and End Semester Examination (ESE). It covers three major areas:

\begin{enumerate}[leftmargin=*]
  \item \textbf{Single-Phase Controlled Rectifiers (5 Hours):} Fully controlled and half-controlled (semi-converter) rectifiers with $RL$ and $RLE$ loads; rectifier and inverter modes of operation; continuous and discontinuous conduction waveforms; output voltage, input line current, real power, power factor ($PF$), and total harmonic distortion ($THD$); effect of source inductance ($L_s$) with commutation overlap angle $\mu$.
  \item \textbf{Three-Phase Controlled Rectifiers (3 Hours):} 6-pulse fully controlled bridge and half-controlled bridge converter with $RLE$ load (continuous conduction, ripple-free current); conduction sequence, firing intervals, waveforms, and output voltage equations; source inductance effect.
  \item \textbf{DC-DC Switching Regulators (4 Hours):} Buck, Boost, and Buck-Boost converters; continuous conduction mode (CCM) operation and waveforms; Inductor Volt-Second Balance and Capacitor Charge Balance; design of filter inductance ($L_{\min}$) and capacitance ($C_{\min}$); qualitative and quantitative effects of non-idealities (inductor resistance $r_L$, capacitor ESR $r_C$); selection criteria for power semiconductor switches (MOSFET, IGBT, Schottky diode).
\end{enumerate}

\begin{conceptbox}[Bloom's Knowledge Level K4 (Analyse)]
In your examinations (Part A 3-mark questions and Part B 9-mark questions), you are expected to:
\begin{itemize}[leftmargin=*]
  \item Sketch clear, aligned waveforms showing supply voltage, gate pulses, load voltage, switch voltage, and input current.
  \item Derive mathematical expressions from first principles using calculus and integral definitions.
  \item Compute exact numerical values for output voltage, RMS quantities, displacement factor, distortion factor, power factor, and filter components.
\end{itemize}
\end{conceptbox}

%%%%%%%%%%%%%%%%%%%%%%%%%%%%%%%%%%%%%%%%%%%%%%%%%%%%%%%%%%%%%%%%%%%%%%%%%%%%%%%
\section{Foundations of Phase Control \& Commutation Mechanics}
%%%%%%%%%%%%%%%%%%%%%%%%%%%%%%%%%%%%%%%%%%%%%%%%%%%%%%%%%%%%%%%%%%%%%%%%%%%%%%%

\subsection{What is Phase-Controlled Rectification?}
An uncontrolled diode conducts whenever its anode-to-cathode voltage is positive ($v_{AK} > 0$). In contrast, a \textbf{Silicon Controlled Rectifier (SCR / Thyristor)} blocks forward voltage until a triggering gate pulse $i_G$ is applied. 
By delaying this gate trigger pulse by an angle $\alpha$ with respect to the AC voltage zero-crossing:
\begin{conceptbox}
\centering\bfseries\color{navyprimary}
Average output DC voltage $V_o$ can be smoothly and steplessly controlled from maximum to zero (or negative).
\end{conceptbox}

\subsection{Natural (Line) Commutation vs. Forced Commutation}
\begin{itemize}[leftmargin=*]
  \item \textbf{Natural (Line) Commutation:} In AC-fed circuits, the AC line voltage reverses polarity every half cycle. When the supply voltage reverses, a negative voltage is placed across the conducting SCR, naturally reducing its current below holding current ($I_H$) and providing reverse bias for turn-off time $t_q$. No auxiliary components are needed.
  \item \textbf{Forced Commutation:} Required when the supply is DC (classical choppers/inverters), necessitating auxiliary capacitors/inductors to force SCR current to zero.
\end{itemize}

\subsection{Key Angular Definitions}
\begin{itemize}[leftmargin=*]
  \item \textbf{Firing Angle (Delay Angle) $\alpha$:} The angle measured from the instant the thyristor would naturally conduct if it were an ordinary diode (the \textit{natural commutation point}) to the instant the gate pulse is injected.
  \item \textbf{Extinction Angle $\beta$:} The instant $\omega t = \beta$ at which the thyristor anode current drops to zero.
  \item \textbf{Conduction Angle $\gamma$:} The angular duration over which the device conducts: $\gamma = \beta - \alpha$.
\end{itemize}

\subsection{Converter Quadrants of Operation}
\begin{conceptbox}[Thyristor Current Unidirectionality]
Because an SCR is a single-direction semiconductor switch, the DC load current $I_o$ can \textbf{never reverse}:
\[
I_o \ge 0 \quad (\text{Strictly unidirectional})
\]
However, the average DC output voltage $V_o$ \textbf{can reverse} polarity depending on $\alpha$ and circuit topology!
\begin{itemize}[leftmargin=*]
  \item \textbf{Quadrant I ($V_o > 0, I_o > 0$):} Rectification Mode. Power flows from AC source to DC load.
  \item \textbf{Quadrant IV ($V_o < 0, I_o > 0$):} Inverter Mode. Power flows from DC active source back into the AC utility grid.
  \item \textbf{Semi-Converter:} Operates in \textbf{Quadrant I only} (One-quadrant converter).
  \item \textbf{Full Converter:} Operates in \textbf{Quadrants I and IV} (Two-quadrant converter).
\end{itemize}
\end{conceptbox}

%%%%%%%%%%%%%%%%%%%%%%%%%%%%%%%%%%%%%%%%%%%%%%%%%%%%%%%%%%%%%%%%%%%%%%%%%%%%%%%
\section{Single-Phase Controlled Rectifiers (Deep Dive)}
%%%%%%%%%%%%%%%%%%%%%%%%%%%%%%%%%%%%%%%%%%%%%%%%%%%%%%%%%%%%%%%%%%%%%%%%%%%%%%%

\subsection{Single-Phase Half-Wave Rectifiers: Complete Load Analysis}

\subsubsection{Resistive Load (\texorpdfstring{$R$}{R})}
\begin{figure}[htbp]
\centering
\includegraphics[width=0.72\textwidth]{figures/1ph_hw_r.png}
\caption{Single-phase half-wave controlled rectifier with resistive load: Circuit and complete waveforms showing firing delay $\alpha$, load voltage $v_o$, and device voltage $v_T$.}
\end{figure}

\textbf{How it works:} The AC supply voltage $v_s = V_m \sin\omega t$ starts its positive half cycle at $\omega t = 0$. Even though the SCR becomes forward-biased at this instant, it \textit{does not conduct} --- it waits in the blocking state until a gate pulse arrives at $\omega t = \alpha$. This delay angle $\alpha$ is our control variable.

Once triggered at $\alpha$, the SCR turns ON and the full supply voltage appears across the resistive load: $v_o = v_s = V_m \sin\omega t$. With a purely resistive load, the current is simply $i_o = v_o / R$, which exactly follows the voltage shape. At $\omega t = \pi$, the supply voltage crosses zero and tries to go negative. Since a resistive load cannot sustain current flow against a negative voltage, the current drops to zero and the SCR naturally turns OFF.

\textbf{Result:} The load sees voltage only during $\alpha \le \omega t \le \pi$ in each cycle, and zero voltage elsewhere.

\textbf{Average (DC) output voltage} is the total area under the $v_o$ curve divided by the full period $2\pi$:
\begin{equation}
  V_o = \frac{1}{2\pi} \int_{\alpha}^{\pi} V_m \sin\omega t \, d(\omega t) = \frac{V_m}{2\pi} [-\cos\omega t]_\alpha^\pi = \frac{V_m}{2\pi} (1 + \cos\alpha)
\end{equation}
\textbf{RMS output voltage} (needed for power and ripple calculations):
\begin{equation}
  V_{o,\text{rms}} = \frac{V_m}{2} \sqrt{\frac{1}{\pi}\left[(\pi - \alpha) + \frac{1}{2}\sin(2\alpha)\right]}
\end{equation}
Notice that at $\alpha = 0$, $V_o = V_m/\pi$ (same as an uncontrolled diode rectifier), and at $\alpha = \pi$, $V_o = 0$ (SCR never conducts).

\subsubsection{Inductive Load (\texorpdfstring{$RL$}{RL}) --- Why Negative Voltage Occurs}
\begin{figure}[htbp]
\centering
\includegraphics[width=0.75\textwidth]{figures/1ph_hw_rl.png}
\caption{Single-phase half-wave circuit with $RL$ load: Waveforms showing conduction extending into the negative half-cycle up to extinction angle $\beta > \pi$ due to inductor induced EMF.}
\end{figure}

When an inductor is present in series with $R$:
\begin{itemize}[leftmargin=*]
  \item During $\alpha \le \omega t \le \pi$, the inductor stores energy in its magnetic field: $E_L = \frac{1}{2} L i^2$.
  \item At $\omega t = \pi$, supply voltage $v_s$ crosses zero and goes negative.
  \item The collapsing magnetic field of the inductor induces a voltage:
  \[
  v_L = L \frac{di}{dt} < 0
  \]
  which acts as a source in series with the AC line, keeping the thyristor anode positive with respect to cathode!
  \item Consequently, the thyristor continues to conduct into the negative half-cycle until all magnetic energy is discharged at extinction angle $\beta > \pi$.
  \item The output voltage is negative during $\pi \le \omega t \le \beta$:
  \[
  V_o = \frac{V_m}{2\pi} (\cos\alpha - \cos\beta)
  \]
\end{itemize}

\begin{formulabox}[General Current Solution for Half-Wave RL Load]
Applying KVL to the circuit during conduction ($\alpha \le \omega t \le \beta$):
\[
V_m \sin\omega t = R\,i + L \frac{di}{d(\omega t)}\,\omega
\]
Let $Z = \sqrt{R^2 + (\omega L)^2}$ be the load impedance and $\phi = \tan^{-1}\!\left(\frac{\omega L}{R}\right)$ the impedance angle. The complete solution (particular + complementary) is:
\[
\boxed{i(\omega t) = \frac{V_m}{Z}\left[\sin(\omega t - \phi) - \sin(\alpha - \phi)\,e^{-\frac{R}{\omega L}(\omega t - \alpha)}\right], \quad \alpha \le \omega t \le \beta}
\]
\begin{itemize}[leftmargin=*]
  \item The first term $\frac{V_m}{Z}\sin(\omega t - \phi)$ is the \textbf{steady-state (forced) response}.
  \item The second term is the \textbf{transient (natural) response} that decays exponentially with time constant $\tau = L/R$.
  \item The extinction angle $\beta$ is found by setting $i(\beta) = 0$ (requires numerical/graphical solution in general).
  \item \textbf{Condition for continuous conduction:} If $\alpha \le \phi$, the thyristor is triggered before the current from the previous cycle has decayed to zero, leading to continuous conduction mode.
\end{itemize}
\end{formulabox}

\subsubsection{Action of Freewheeling Diode (\texorpdfstring{$D_{\text{fw}}$}{Dfw})}
\begin{figure}[htbp]
\centering
\includegraphics[width=0.72\textwidth]{figures/1ph_hw_rl_fwd.png}
\caption{Single-phase half-wave circuit with $RL$ load and Freewheeling Diode (FWD): Negative voltage is completely eliminated as current freewheels through $D_{fw}$ from $\pi$ to $\beta$.}
\end{figure}

Connecting an anti-parallel diode $D_{fw}$ across the $RL$ load completely eliminates the negative voltage excursion:
\begin{itemize}[leftmargin=*]
  \item At $\omega t = \pi$, as $v_s$ goes negative, the cathode of $D_{fw}$ becomes negative relative to its anode. $D_{fw}$ turns ON automatically.
  \item $D_{fw}$ clamps the load voltage to zero: $v_o(t) \approx 0$ for $\pi \le \omega t \le \beta$.
  \item The load current freewheels through $R, L,$ and $D_{fw}$, decaying exponentially.
  \item The reverse voltage $v_s$ immediately appears across the SCR, commutating it OFF at $\omega t = \pi$.
\end{itemize}

\subsubsection{Active Load (\texorpdfstring{$RLE$}{RLE} / Battery Charging)}
\begin{figure}[htbp]
\centering
\includegraphics[width=0.72\textwidth]{figures/1ph_hw_rle.png}
\caption{Single-phase half-wave circuit with $RLE$ load: Thyristor is forward biased only when $v_s > E$. Firing angle $\alpha$ cannot be less than $\theta_1 = \sin^{-1}(E/V_m)$.}
\end{figure}

An $RLE$ load models a practical DC motor drive or battery charger: $R$ represents the winding/internal resistance, $L$ is the motor armature inductance or smoothing inductor, and $E$ is the motor back-EMF or battery terminal voltage.

\textbf{Key difference from $R$ and $RL$ loads:} The thyristor cannot conduct unless the supply voltage exceeds the back-EMF, i.e., $v_s > E$. This creates a \textit{minimum firing angle} constraint:
\[
\theta_1 = \sin^{-1}\!\left(\frac{E}{V_m}\right) \quad \implies \quad \alpha \ge \theta_1
\]
If we try to fire the SCR at $\alpha < \theta_1$, the net voltage $(v_s - E)$ across $R$ and $L$ is negative --- no current can flow, and the gate pulse is wasted.

Once fired at $\alpha \ge \theta_1$, the load current builds up according to: $v_s - E = R\,i + L\,di/dt$. Current flows until the stored magnetic energy in $L$ is fully dissipated (at extinction angle $\beta$). The current waveform shape depends on the $L/R$ time constant and the magnitude of $E$.

\subsection{Single-Phase Fully Controlled Bridge Rectifier (Full Converter)}

A full converter uses \textbf{four thyristors} arranged in a bridge: $T_1, T_2$ in one diagonal and $T_3, T_4$ in the other. Unlike the half-wave circuit which uses only one half of each AC cycle, the full bridge rectifies \textit{both} halves, doubling the average DC output.

\subsubsection{Pure Resistive Load (\texorpdfstring{$R$}{R})}
\begin{figure}[htbp]
\centering
\includegraphics[width=0.48\textwidth]{figures/1ph_full_r_ckt.png}
\hfill
\includegraphics[width=0.48\textwidth]{figures/1ph_full_r_wave.png}
\caption{Single-phase fully controlled bridge with resistive load: Circuit schematic and synchronized voltage and current waveforms.}
\end{figure}

\textbf{How the bridge works with R load:}
\begin{itemize}[leftmargin=*]
  \item \textbf{Positive half cycle ($\alpha \le \omega t \le \pi$):} Thyristors $T_1$ and $T_2$ are fired simultaneously. Current flows: source $\to T_1 \to R \to T_2 \to$ source. Load voltage $v_o = v_s > 0$.
  \item At $\omega t = \pi$, supply reverses. With a resistive load, current drops to zero and $T_1, T_2$ turn OFF naturally.
  \item \textbf{Negative half cycle ($\pi+\alpha \le \omega t \le 2\pi$):} Thyristors $T_3$ and $T_4$ are fired. Current flows in the opposite path through the source but in the \textit{same direction} through the load. Load voltage $v_o = -v_s > 0$.
  \item At $2\pi$, supply reverses again, current drops to zero, and the cycle repeats.
\end{itemize}
The result is two positive voltage pulses per cycle, each of width $(\pi - \alpha)$. This is why the full converter produces twice the average voltage of a half-wave rectifier.

\subsubsection{Continuous Ripple-Free Conduction (\texorpdfstring{$RL$ / $RLE$ Load}{RL / RLE Load})}
\begin{figure}[htbp]
\centering
\includegraphics[width=0.48\textwidth]{figures/1ph_full_rl_ckt.png}
\hfill
\includegraphics[width=0.48\textwidth]{figures/1ph_full_rl_wave.png}
\caption{Single-phase fully controlled bridge with highly inductive load: Circuit schematic and continuous ripple-free waveforms showing negative voltage loops.}
\end{figure}

\begin{figure}[htbp]
\centering
\includegraphics[width=0.75\textwidth]{figures/1ph_full_rl_wave2.png}
\caption{Single-phase Full Converter detailed waveforms: Input voltage $v_s$, gate pulses $i_{G1,2}$ and $i_{G3,4}$, output voltage $v_o$, constant load current $I_o$, input line current $i_s$, thyristor currents, and device voltage $v_{T1}$.}
\end{figure}

Assuming load inductance is very large ($\omega L \gg R$), the load current is a smooth, continuous DC value $I_o$. This is the most important case for analysis.

\textbf{The key idea:} Unlike a resistive load (where current drops to zero when voltage reverses), a large inductor \textit{forces the current to continue flowing} even when the supply voltage goes negative. The inductor ``fights'' any change in current --- it simply will not allow $I_o$ to drop. This creates negative voltage intervals across the load, which is the defining feature of a full converter with inductive load.

\begin{itemize}[leftmargin=*]
  \item \textbf{Interval 1 ($\alpha \le \omega t \le \pi$) --- Normal Rectification:} Thyristors $T_1, T_2$ are fired at $\omega t = \alpha$. The supply voltage is positive, so $v_o = v_s > 0$. Current $I_o$ flows through $T_1, T_2$. The inductor stores energy. Source current $i_s = +I_o$.
  \item \textbf{Interval 2 ($\pi \le \omega t \le \pi+\alpha$) --- Inductor Holds On:} At $\omega t = \pi$, the supply voltage reverses and becomes negative. \textit{But} $T_1, T_2$ cannot turn off --- the enormous inductor maintains its current and keeps $T_1, T_2$ forward-biased. The load voltage goes negative: $v_o = v_s < 0$. During this interval, energy stored in the inductor is being returned to the AC source. Think of it as the inductor ``pushing current uphill'' against the reversed supply.
  \item \textbf{Interval 3 ($\pi+\alpha \le \omega t \le 2\pi$) --- Commutation:} At $\omega t = \pi+\alpha$, thyristors $T_3, T_4$ receive their gate pulses and turn ON. Now the line voltage instantly reverse-biases $T_1, T_2$ (because $T_3, T_4$ connect the source in opposite polarity across $T_1, T_2$), forcing them OFF. This is natural (line) commutation. Now $T_3, T_4$ carry the load current, and $v_o = -v_s > 0$. Source current reverses: $i_s = -I_o$.
  \item \textbf{Interval 4 ($2\pi \le \omega t \le 2\pi+\alpha$) --- Mirror of Interval 2:} $T_3, T_4$ continue conducting through the next zero crossing. Load voltage is again negative as the inductor holds conduction.
\end{itemize}

\subsubsection{Derivation of Average Output Voltage (\texorpdfstring{$V_o$}{Vo})}
\textbf{Why integrate from $\alpha$ to $\pi+\alpha$?} Look at the waveforms above. The output voltage $v_o$ repeats every $\pi$ radians (not $2\pi$), because the full bridge rectifies both halves of the AC cycle. One complete output period runs from $\omega t = \alpha$ (when a thyristor pair is fired) to $\omega t = \pi+\alpha$ (when the next pair takes over). So the average is:
\[
V_o = \frac{1}{\pi} \int_{\alpha}^{\pi+\alpha} V_m \sin\omega t \, d(\omega t) = \frac{V_m}{\pi} [-\cos\omega t]_\alpha^{\pi+\alpha}
\]
\[
V_o = \frac{V_m}{\pi} [-\cos(\pi+\alpha) - (-\cos\alpha)] = \frac{V_m}{\pi} [\cos\alpha + \cos\alpha]
\]
\begin{formulabox}[Single-Phase Full Converter Output Voltage]
\[
\boxed{V_o = \frac{2 V_m}{\pi} \cos\alpha = V_{o0} \cos\alpha}
\]
where $V_{o0} = \frac{2 V_m}{\pi} \approx 0.6366 V_m$ is the maximum average output voltage at $\alpha = 0^\circ$.

\textbf{Physical insight:} $V_o$ is positive for $\alpha < 90^\circ$ (rectification), zero at $\alpha = 90^\circ$, and negative for $\alpha > 90^\circ$ (inversion). The $\cos\alpha$ dependence means voltage control is non-linear --- most sensitive near $\alpha = 90^\circ$.
\end{formulabox}

\subsubsection{Derivation of RMS Output Voltage (\texorpdfstring{$V_{o,\text{rms}}$}{Vo,rms})}
\[
V_{o,\text{rms}} = \sqrt{\frac{1}{\pi} \int_{\alpha}^{\pi+\alpha} V_m^2 \sin^2\omega t \, d(\omega t)} = \frac{V_m}{\sqrt{2}} = V_s
\]
\begin{conceptbox}[RMS Voltage Independence]
For continuous conduction in a full converter, \textbf{$V_{o,\text{rms}}$ is always equal to the supply RMS voltage $V_s$}, regardless of firing angle $\alpha$! This is because the integration of $\sin^2$ over any complete half-cycle always gives the same result. This means the output AC ripple power increases as $\alpha$ increases (since average DC power decreases but RMS stays constant).
\end{conceptbox}

\subsubsection{Inverter Mode of Operation (\texorpdfstring{$\alpha > 90^\circ$}{alpha > 90 deg})}
\begin{figure}[htbp]
\centering
\includegraphics[width=0.72\textwidth]{figures/1ph_full_inverter.png}
\caption{Line-Commutated Inverter Operation: When $\alpha > 90^\circ$, average output voltage $V_o$ becomes negative. An active DC source $E$ connected in reverse polarity drives current through the thyristors, regenerating power into the AC line.}
\end{figure}

\textbf{When would you want negative output voltage?} Consider a DC motor decelerating: its back-EMF $E$ acts as a source. By increasing $\alpha$ beyond $90^\circ$, the converter output voltage reverses, allowing the motor's kinetic energy to flow \textit{back} into the AC grid. This is called \textbf{regenerative braking}. The current direction stays the same (SCRs are unidirectional), but the voltage reverses --- so power $P = V_o \times I_o$ becomes negative, meaning the load is returning power.

\begin{figure}[htbp]
\centering
\includegraphics[width=0.75\textwidth]{figures/1ph_full_rle_wave.png}
\caption{Single-phase full converter with $RLE$ load: Detailed voltage and current waveforms in rectifying mode ($\alpha < 90^\circ$) and inverting mode ($\alpha > 90^\circ$).}
\end{figure}

\begin{alertbox}[Commutation Failure in Inverter Mode]
During inversion ($\alpha > 90^\circ$), the reverse voltage available across the commutating thyristor exists only for the margin angle:
\[
\gamma = 180^\circ - \alpha
\]
The thyristor requires a physical turn-off time $t_q$. If $\gamma < \omega t_q$, the outgoing thyristor fails to recover its forward blocking capability before $v_s$ becomes positive. Both bridge legs conduct simultaneously, causing a massive short circuit known as \textbf{Commutation Failure}. To provide safety against this, practical firing angles are restricted to:
\[
\alpha_{\max} \approx 165^\circ \text{ to } 170^\circ
\]
\end{alertbox}

\subsection{AC Source Performance Metrics (Fourier Analysis of Input Current)}

\textbf{Why is the input current a square wave?} Recall from the interval analysis above: during $\alpha \le \omega t \le \pi+\alpha$, thyristors $T_1, T_2$ conduct and the source current is $i_s = +I_o$ (constant, because load current is ripple-free). Then during $\pi+\alpha \le \omega t \le 2\pi+\alpha$, thyristors $T_3, T_4$ conduct and the source current flips to $i_s = -I_o$. So the input current alternates between $+I_o$ and $-I_o$ --- a perfect \textbf{square wave} of amplitude $I_o$, delayed by angle $\alpha$ from the supply voltage zero crossing.

\subsubsection{Fourier Series Expansion}
Since the square wave has odd symmetry and half-wave symmetry, only odd harmonics exist:
\[
i_s(t) = \sum_{n=1,3,5,\dots}^{\infty} \frac{4 I_o}{n \pi} \sin(n(\omega t - \alpha))
\]
\begin{itemize}[leftmargin=*]
  \item \textbf{Fundamental Peak Current:} $I_{m1} = \frac{4 I_o}{\pi}$
  \item \textbf{Fundamental RMS Current (\texorpdfstring{$I_{s1}$}{Is1}):}
  \[
  \boxed{I_{s1} = \frac{4 I_o}{\sqrt{2} \pi} = \frac{2\sqrt{2}}{\pi} I_o \approx 0.9003 I_o}
  \]
  \item \textbf{Total RMS Input Current (\texorpdfstring{$I_s$}{Is}):}
  \[
  \boxed{I_s = \sqrt{\frac{1}{\pi} \int_{\alpha}^{\pi+\alpha} I_o^2 \, d(\omega t)} = I_o}
  \]
\end{itemize}

\subsubsection{Power Factor Decomposition}
Active power transmitted to the DC load is:
\[
P = V_s I_{s1} \cos\phi_1 = V_s \left(\frac{2\sqrt{2}}{\pi} I_o\right) \cos\alpha = V_o I_o
\]
where $\phi_1 = \alpha$ is the fundamental displacement angle.
\begin{formulabox}[Harmonic Performance Indices]
\begin{enumerate}[leftmargin=*]
  \item \textbf{Displacement Factor ($DF$):} $DF = \cos\phi_1 = \cos\alpha$
  \item \textbf{Current Distortion Factor ($g$):} $g = \frac{I_{s1}}{I_s} = \frac{\frac{2\sqrt{2}}{\pi} I_o}{I_o} = \frac{2\sqrt{2}}{\pi} \approx 0.9003$
  \item \textbf{Input Power Factor ($PF$):} $PF = \frac{P}{S} = g \times DF = \boxed{\frac{2\sqrt{2}}{\pi} \cos\alpha \approx 0.9003 \cos\alpha}$
  \item \textbf{Total Harmonic Distortion ($THD$):} $THD = \sqrt{\frac{I_s^2 - I_{s1}^2}{I_{s1}^2}} = \sqrt{\frac{1}{g^2} - 1} = \sqrt{\frac{\pi^2}{8} - 1} \approx \boxed{48.34\%}$
\end{enumerate}
\end{formulabox}

\subsection{Single-Phase Half-Controlled Converter (Semi-Converter)}

\textbf{What is a Semi-Converter and why use it?}
A full converter uses four thyristors, requiring four separate gate trigger channels with pulse transformers or optocouplers. However, many practical industrial DC loads (like battery chargers, electroplating, and unidirectional DC motor drives) only operate in **Quadrant I** ($V_o > 0, I_o > 0$) and never need regenerative inversion. 

In such applications, replacing two thyristors with passive uncontrolled power diodes creates a \textbf{Semi-Converter (Half-Controlled Bridge)}. This provides several major engineering advantages:
\begin{enumerate}[leftmargin=*]
  \item \textbf{Lower Cost \& Simpler Controls:} Diodes are cheaper than SCRs and require no gate drive circuits, power supplies, or synchronisation logic.
  \item \textbf{Built-in Freewheeling Action:} The diodes inherently provide a freewheeling path across the load whenever the AC voltage reverses, completely eliminating negative output voltage loops without needing a separate external freewheeling diode.
  \item \textbf{Superior Power Factor \& Reduced Reactive Power:} Because load current freewheels during $\pi \le \omega t \le \pi+\alpha$, the AC supply draws zero current during this interval ($i_s = 0$). This reduces the reactive power drawn from the grid and gives a significantly higher power factor than a full converter at all firing angles $\alpha > 0$.
\end{enumerate}

\begin{figure}[htbp]
\centering
\includegraphics[width=0.48\textwidth]{figures/1ph_semi_symm.png}
\hfill
\includegraphics[width=0.48\textwidth]{figures/1ph_semi_asymm.png}
\caption{Single-phase Semi-Converter configurations: (Left) Symmetrical bridge with thyristors in top bus and diodes in bottom bus. (Right) Asymmetrical bridge with one thyristor and one diode per leg.}
\end{figure}

\subsubsection{Symmetrical vs. Asymmetrical Configurations}
\begin{itemize}[leftmargin=*]
  \item \textbf{Symmetrical Bridge:} Thyristors $T_1, T_2$ are placed in the positive (top) arm, and diodes $D_1, D_2$ in the negative (bottom) arm. During freewheeling, current circulates through one thyristor and one diode (e.g., $T_1$ and $D_2$). Therefore, the thyristor must carry current during the freewheeling period, which prolongs its conduction angle to $(\pi+\alpha)$ and gives a circuit turn-off time of $t_c = (\pi - \alpha)/\omega$.
  \item \textbf{Asymmetrical Bridge:} Each of the two bridge legs contains one thyristor and one diode ($T_1, D_1$ in leg 1, and $T_2, D_2$ in leg 2). Here, freewheeling occurs strictly through the two series diodes ($D_1$ and $D_2$). Both thyristors turn off cleanly at $\omega t = \pi$, providing a full half-cycle turn-off margin $t_c = \pi/\omega$, which makes commutation much more robust.
\end{itemize}

\begin{figure}[htbp]
\centering
\includegraphics[width=0.48\textwidth]{figures/1ph_semi_wave_cont.png}
\hfill
\includegraphics[width=0.48\textwidth]{figures/1ph_semi_wave_discont.png}
\caption{Semi-Converter Waveforms: (Left) Continuous conduction mode showing voltage clamped to zero during freewheeling intervals ($\pi$ to $\pi+\alpha$). (Right) Discontinuous conduction mode.}
\end{figure}

\begin{figure}[htbp]
\centering
\includegraphics[width=0.55\textwidth]{figures/1ph_vo_vs_alpha.png}
\caption{Normalized Average Output Voltage ($V_o / V_{o0}$) vs Firing Angle ($\alpha$): Comparison showing that the full converter can produce negative voltage (inverting mode), while the semi-converter is strictly non-negative ($V_o \ge 0$).}
\end{figure}

\subsubsection{Operating Intervals \& Mathematical Derivation of \texorpdfstring{$V_o$}{Vo}}
Assuming continuous, ripple-free load current $I_o$:
\begin{itemize}[leftmargin=*]
  \item \textbf{Interval 1 ($\alpha \le \omega t \le \pi$) --- Power Delivery:} Thyristor $T_1$ and diode $D_1$ conduct. The load is connected directly to the AC source: $v_o(t) = v_s(t) = V_m \sin\omega t$. Active power flows from the AC utility grid into the DC load. Source current is $i_s = +I_o$.
  \item \textbf{Interval 2 ($\pi \le \omega t \le \pi+\alpha$) --- Freewheeling Mode:} At $\omega t = \pi$, supply voltage $v_s$ crosses zero and turns negative. Diode $D_2$ immediately becomes forward-biased by the collapsing inductor EMF. Current shifts from the AC source into the freewheeling path ($D_1, D_2$). The load terminals are short-circuited through the conducting diodes, clamping the load voltage to zero: $v_o(t) = 0$. During this interval, the AC source is disconnected: $i_s(t) = 0$.
\end{itemize}

\textbf{Average Output Voltage Derivation:}
Since the output voltage is $V_m\sin\omega t$ from $\alpha$ to $\pi$ and $0$ from $\pi$ to $\pi+\alpha$, we integrate over the repeating half-cycle period $\pi$:
\[
V_o = \frac{1}{\pi} \left[ \int_{\alpha}^{\pi} V_m \sin\omega t \, d(\omega t) + \int_{\pi}^{\pi+\alpha} 0 \, d(\omega t) \right] = \frac{V_m}{\pi} [-\cos\omega t]_\alpha^\pi
\]
\[
V_o = \frac{V_m}{\pi} [-\cos\pi - (-\cos\alpha)] = \frac{V_m}{\pi} [1 + \cos\alpha]
\]
\begin{formulabox}[Single-Phase Semi-Converter Output Voltage]
\[
\boxed{V_o = \frac{V_m}{\pi} (1 + \cos\alpha) = \frac{\sqrt{2}V_s}{\pi} (1 + \cos\alpha)}
\]
\begin{itemize}[leftmargin=*]
  \item At $\alpha = 0^\circ$: $V_o = \frac{2 V_m}{\pi}$ (same as full converter at $\alpha = 0$).
  \item At $\alpha = 90^\circ$: $V_o = \frac{V_m}{\pi}$ (half the maximum DC voltage).
  \item At $\alpha = 180^\circ$: $V_o = 0$. Because $(1+\cos\alpha) \ge 0$ for all $\alpha \in [0, \pi]$, $V_o$ can \textbf{never be negative}!
\end{itemize}
\end{formulabox}

\subsubsection{AC Source Current \& Power Factor Advantage}
The source current $i_s(t)$ flows only during the active conduction intervals ($\alpha$ to $\pi$, and $\pi+\alpha$ to $2\pi$), each of angular width $(\pi - \alpha)$:
\begin{itemize}[leftmargin=*]
  \item \textbf{RMS Supply Current:}
  \[
  I_s = \sqrt{\frac{1}{\pi} \int_{\alpha}^{\pi} I_o^2 \, d(\omega t)} = \boxed{I_o \sqrt{\frac{\pi - \alpha}{\pi}}}
  \]
  Notice that as $\alpha$ increases, $I_s$ drops because the current freewheels for longer periods.
  \item \textbf{Displacement Factor ($DF$):} The fundamental input current is centered in the conduction interval $[\alpha, \pi]$, meaning its phase shift relative to supply voltage is only $\phi_1 = \alpha/2$:
  \[
  \boxed{DF = \cos\left(\frac{\alpha}{2}\right)}
  \]
  Compare this to a full converter where $DF = \cos\alpha$. For example, at $\alpha = 60^\circ$, full converter $DF = 0.50$, whereas semi-converter $DF = \cos 30^\circ = 0.866$ --- a 73\% improvement!
  \item \textbf{Input Power Factor ($PF$):}
  \[
  PF = \frac{P}{S} = \frac{V_o I_o}{V_s I_s} = \frac{\frac{V_m}{\pi}(1+\cos\alpha) I_o}{V_s \left(I_o \sqrt{\frac{\pi-\alpha}{\pi}}\right)} = \boxed{\frac{\sqrt{2}(1 + \cos\alpha)}{\sqrt{\pi(\pi - \alpha)}}}
  \]
\end{itemize}

%%%%%%%%%%%%%%%%%%%%%%%%%%%%%%%%%%%%%%%%%%%%%%%%%%%%%%%%%%%%%%%%%%%%%%%%%%%%%%%
\section{Effect of Source Inductance (\texorpdfstring{$L_s$}{Ls}) on Controlled Rectifiers}
%%%%%%%%%%%%%%%%%%%%%%%%%%%%%%%%%%%%%%%%%%%%%%%%%%%%%%%%%%%%%%%%%%%%%%%%%%%%%%%

In practical installations, the AC source possesses leakage inductance from feeding transformers and supply distribution lines ($L_s$).

\subsection{Physical Mechanism of Commutation Overlap}
\begin{figure}[htbp]
\centering
\includegraphics[width=0.48\textwidth]{figures/source_ind_p1.png}
\hfill
\includegraphics[width=0.48\textwidth]{figures/source_ind_p2.png}
\caption{Effect of Source Inductance ($L_s$): (Left) Circuit schematic showing supply inductance $L_s$. (Right) Commutation overlap angle $\mu$ during which all 4 thyristors conduct simultaneously and output voltage is shorted ($v_o = 0$).}
\end{figure}

Because of $L_s$, the line current cannot change instantaneously:
\[
v_{L_s} = L_s \frac{di_s}{dt}
\]
When incoming thyristors $T_1, T_2$ are triggered at $\omega t = \alpha$, outgoing thyristors $T_3, T_4$ cannot shut off abruptly. For an angular duration $\mu$ (the \textbf{overlap angle}), \textbf{all four thyristors conduct simultaneously}.
During this interval, the AC source is shorted through $L_s$ and the bridge switches, causing the output voltage to be \textbf{zero}:
\[
v_o(t) = 0 \quad \text{for } \alpha \le \omega t \le \alpha + \mu
\]

\subsection{Mathematical Derivation of Overlap Angle \texorpdfstring{$\mu$}{mu}}
During overlap ($\alpha \le \omega t \le \alpha+\mu$):
Let $i_1$ be incoming SCR current and $i_2$ be outgoing SCR current:
\[
i_1 + i_2 = I_o \implies \frac{di_2}{dt} = -\frac{di_1}{dt}
\]
Line current is $i_s = i_1 - i_2$. Therefore:
\[
\frac{di_s}{dt} = \frac{di_1}{dt} - \frac{di_2}{dt} = 2 \frac{di_1}{dt}
\]
KVL around the AC source loop gives:
\[
v_s(t) = L_s \frac{di_s}{dt} = 2 L_s \frac{di_1}{dt} \implies \frac{di_1}{d(\omega t)} = \frac{V_m}{2\omega L_s} \sin\omega t
\]
Integrating from $\omega t = \alpha$ ($i_1 = 0$) to $\omega t = \alpha+\mu$ ($i_1 = I_o$):
\[
\int_{0}^{I_o} di_1 = \frac{V_m}{2\omega L_s} \int_{\alpha}^{\alpha+\mu} \sin\omega t \, d(\omega t)
\]
\[
I_o = \frac{V_m}{2\omega L_s} [\cos\alpha - \cos(\alpha+\mu)]
\]
\begin{formulabox}[Commutation Overlap Equation]
\[
\boxed{\cos\alpha - \cos(\alpha+\mu) = \frac{2\omega L_s I_o}{V_m} = \frac{\sqrt{2}\omega L_s I_o}{V_s}}
\]
\end{formulabox}

\subsection{Average Output Voltage Reduction \& Equivalent Commutation Resistance}
\begin{figure}[htbp]
\centering
\includegraphics[width=0.55\textwidth]{figures/source_ind_p6.png}
\caption{Detailed voltage and current commutation waveforms: Commutation interval $\mu$, decay of outgoing current $i_2$, rise of incoming current $i_1$, and the resulting lost volt-second area $\Delta V_o$.}
\end{figure}

\begin{figure}[htbp]
\centering
\includegraphics[width=0.48\textwidth]{figures/source_ind_p7.png}
\caption{DC Equivalent Circuit with Source Inductance: Converter modeled as ideal DC voltage source $V_{o0}\cos\alpha$ in series with equivalent commutation resistance $R_c = 2 f L_s$.}
\end{figure}

Since $v_o(t) = 0$ during overlap, conduction only provides voltage from $\alpha+\mu$ to $\pi+\alpha$:
\[
V_o = \frac{1}{\pi} \int_{\alpha+\mu}^{\pi+\alpha} V_m \sin\omega t \, d(\omega t) = \frac{V_m}{\pi} [\cos\alpha + \cos(\alpha+\mu)]
\]
Substituting $\cos(\alpha+\mu) = \cos\alpha - \frac{2\omega L_s I_o}{V_m}$:
\[
V_o = \frac{V_m}{\pi} \left[ 2\cos\alpha - \frac{2\omega L_s I_o}{V_m} \right] = \frac{2 V_m}{\pi}\cos\alpha - \frac{\omega L_s}{\pi} I_o
\]
\begin{formulabox}[Output Voltage with Source Inductance]
\[
\boxed{V_o = \frac{2 V_m}{\pi}\cos\alpha - \Delta V_o = V_{o0}\cos\alpha - R_c I_o}
\]
where:
\[
\Delta V_o = \frac{\omega L_s}{\pi} I_o = 2 f L_s I_o, \qquad \boxed{R_c = \frac{\omega L_s}{\pi} = 2 f L_s \quad (\Omega)}
\]
\end{formulabox}

\begin{conceptbox}[Nature of Equivalent Commutation Resistance $R_c$]
$R_c$ is an \textbf{apparent} or \textbf{fictitious resistance}. It represents the loss of average voltage area due to overlap. \textbf{No thermal heat loss ($I_o^2 R_c$) occurs in $R_c$}!
\end{conceptbox}

%%%%%%%%%%%%%%%%%%%%%%%%%%%%%%%%%%%%%%%%%%%%%%%%%%%%%%%%%%%%%%%%%%%%%%%%%%%%%%%
\section{Three-Phase Controlled Rectifiers}
%%%%%%%%%%%%%%%%%%%%%%%%%%%%%%%%%%%%%%%%%%%%%%%%%%%%%%%%%%%%%%%%%%%%%%%%%%%%%%%

Three-phase rectifiers provide higher DC power capability, lower ripple factor, and a ripple frequency of $6 f_s = 300\text{ Hz}$ (for $50\text{ Hz}$ supply).

\begin{conceptbox}[Three-Phase Voltage Relationships --- Quick Reference]
For a balanced three-phase supply with RMS line-to-line voltage $V_L$ and RMS phase voltage $V_{ph}$:
\[
\underbrace{V_{mp}}_{\text{Peak Phase}} = \sqrt{2}\,V_{ph}, \qquad V_L = \sqrt{3}\,V_{ph}, \qquad \underbrace{V_{ml}}_{\text{Peak Line}} = \sqrt{2}\,V_L = \sqrt{3}\,V_{mp}
\]
\textbf{Example:} For $V_L = 400\text{ V}$: $V_{ph} = 231\text{ V}$, $V_{mp} = 326.6\text{ V}$, $V_{ml} = 565.7\text{ V}$.
\end{conceptbox}

\subsection{Three-Phase Half-Wave Converter (3-Pulse Rectifier)}
\begin{figure}[htbp]
\centering
\includegraphics[width=0.48\textwidth]{figures/3ph_hw_ckt.png}
\hfill
\includegraphics[width=0.48\textwidth]{figures/3ph_hw_wave_a0.png}
\caption{Three-phase half-wave controlled rectifier: (Left) Circuit schematic with 3 thyristors. (Right) Output voltage waveform for $\alpha = 0^\circ$.}
\end{figure}

\begin{figure}[htbp]
\centering
\includegraphics[width=0.48\textwidth]{figures/3ph_hw_wave_a30.png}
\hfill
\includegraphics[width=0.48\textwidth]{figures/3ph_hw_wave_a60.png}
\caption{Three-phase half-wave waveforms at delayed firing: (Left) $\alpha = 30^\circ$. (Right) $\alpha = 60^\circ$.}
\end{figure}

Each thyristor conducts for $120^\circ$ ($2\pi/3$ radians). There are 3 pulses per cycle of period $2\pi/3$. Integrating the phase voltage over one conduction interval:
\[
V_o = \frac{1}{2\pi/3} \int_{-\pi/6+\alpha}^{\pi/2+\alpha} V_{mp} \sin\omega t \, d(\omega t) = \frac{3 V_{mp}}{2\pi} [-\cos\omega t]_{-\pi/6+\alpha}^{\pi/2+\alpha}
\]
Using $\cos(-\pi/6+\alpha) - \cos(\pi/2+\alpha) = 2\cos\alpha \sin(60^\circ) = \sqrt{3}\cos\alpha$:
\begin{formulabox}[3-Phase Half-Wave (3-Pulse) Output Voltage]
\[
\boxed{V_o = \frac{3\sqrt{3}\,V_{mp}}{2\pi}\cos\alpha = \frac{3\,V_{ml}}{2\pi}\cos\alpha \approx 1.17\,V_{ph}\cos\alpha}
\]
where $V_{mp}$ is the peak phase voltage and $V_{ml} = \sqrt{3}\,V_{mp}$ is the peak line-to-line voltage.
\end{formulabox}

\subsection{Three-Phase Fully Controlled 6-Pulse Bridge Converter}
Contains 6 thyristors:
\begin{itemize}[leftmargin=*]
  \item \textbf{Positive Group (Cathodes common to DC+ bus):} $T_1$ (phase $a$), $T_3$ (phase $b$), $T_5$ (phase $c$).
  \item \textbf{Negative Group (Anodes common to DC$-$ bus):} $T_4$ (phase $a$), $T_6$ (phase $b$), $T_2$ (phase $c$).
\end{itemize}

\begin{figure}[htbp]
\centering
\includegraphics[width=0.55\textwidth]{figures/3ph_full_ckt.png}
\caption{Three-Phase Fully Controlled Bridge Converter (6-Pulse): Schematic diagram showing positive group ($T_1, T_3, T_5$) and negative group ($T_4, T_6, T_2$).}
\end{figure}

\subsubsection{Firing Sequence and Conduction Table}
Thyristors are triggered in numerical order at $60^\circ$ intervals:
\[
\boxed{T_1 \longrightarrow T_2 \longrightarrow T_3 \longrightarrow T_4 \longrightarrow T_5 \longrightarrow T_6 \longrightarrow T_1}
\]
Each SCR conducts for $120^\circ$. At any instant, two devices conduct (one positive group, one negative group).

\begin{table}[h!]
\centering
\small
\caption{Conduction Intervals and Line Voltages for 3-Phase Full Converter}
\begin{tabular}{ccc}
\toprule
\textbf{Interval ($\omega t$)} & \textbf{Conducting Pair} & \textbf{Output Voltage $v_o(t)$} \\
\midrule
$60^\circ + \alpha$ to $120^\circ + \alpha$ & $T_1, T_6$ & $v_{ab} = v_{an} - v_{bn}$ \\
$120^\circ + \alpha$ to $180^\circ + \alpha$ & $T_1, T_2$ & $v_{ac} = v_{an} - v_{cn}$ \\
$180^\circ + \alpha$ to $240^\circ + \alpha$ & $T_3, T_2$ & $v_{bc} = v_{bn} - v_{cn}$ \\
$240^\circ + \alpha$ to $300^\circ + \alpha$ & $T_3, T_4$ & $v_{ba} = v_{bn} - v_{an}$ \\
$300^\circ + \alpha$ to $360^\circ + \alpha$ & $T_5, T_4$ & $v_{ca} = v_{cn} - v_{an}$ \\
$360^\circ + \alpha$ to $60^\circ + \alpha$ & $T_5, T_6$ & $v_{cb} = v_{cn} - v_{bn}$ \\
\bottomrule
\end{tabular}
\end{table}

\subsubsection{Waveform Evolution with Firing Angle $\alpha$}
\begin{figure}[htbp]
\centering
\includegraphics[width=0.48\textwidth]{figures/3ph_full_wave_a0.png}
\hfill
\includegraphics[width=0.48\textwidth]{figures/3ph_full_wave_a30.png}
\caption{Three-phase Full Converter Output Voltage Waveforms: (Left) $\alpha = 0^\circ$ (pure 6-pulse diode envelope). (Right) $\alpha = 30^\circ$.}
\end{figure}

\begin{figure}[htbp]
\centering
\includegraphics[width=0.48\textwidth]{figures/3ph_full_wave_a60.png}
\hfill
\includegraphics[width=0.48\textwidth]{figures/3ph_full_wave_a90.png}
\caption{Three-phase Full Converter Waveforms: (Left) $\alpha = 60^\circ$ (instantaneous voltage touches zero). (Right) $\alpha = 90^\circ$ (voltage swings equally positive and negative, average $V_o = 0$).}
\end{figure}

\subsubsection{Derivation of Average Output Voltage (\texorpdfstring{$V_o$}{Vo})}
There are 6 symmetric pulses per cycle of duration $\pi/3$ ($60^\circ$):
\[
V_o = \frac{1}{\pi/3} \int_{-\pi/6+\alpha}^{\pi/6+\alpha} V_{ml} \cos\omega t \, d(\omega t) = \frac{3 V_{ml}}{\pi} [\sin\omega t]_{-\pi/6+\alpha}^{\pi/6+\alpha}
\]
Using $\sin(A+B) - \sin(A-B) = 2 \cos A \sin B$:
\[
V_o = \frac{3 V_{ml}}{\pi} \left[ 2 \cos\alpha \sin(30^\circ) \right]
\]
\begin{formulabox}[3-Phase Full Converter Output Voltage]
\[
\boxed{V_o = \frac{3 V_{ml}}{\pi} \cos\alpha = \frac{3\sqrt{3} V_{mp}}{\pi} \cos\alpha = \frac{3\sqrt{2} V_L}{\pi} \cos\alpha \approx 1.35 V_L \cos\alpha}
\]
\end{formulabox}

\subsection{Three-Phase Half-Controlled Converter (Semi-Converter)}

\textbf{Circuit Topology \& Operating Philosophy:}
The 3-phase semi-converter replaces the 3 thyristors in the negative group with 3 power diodes ($D_4, D_6, D_2$), while retaining 3 thyristors ($T_1, T_3, T_5$) in the positive group (or vice versa). 
\begin{itemize}[leftmargin=*]
  \item \textbf{Why use it?} It reduces the number of gate firing circuits from 6 down to 3, significantly lowering control complexity and hardware cost.
  \item \textbf{Single-Quadrant Operation:} Because of the passive diodes, the DC terminal voltage $V_o$ can \textbf{never reverse polarity} ($V_o \ge 0$). Thus, it operates strictly in \textbf{Quadrant I} as a controlled rectifier.
  \item \textbf{Higher Power Factor:} Just like the 1-phase semi-converter, the diodes provide an automatic freewheeling action that clamps negative voltage excursions, eliminating reactive power circulation and yielding a higher power factor than the 6-pulse full bridge for $\alpha > 0$.
\end{itemize}

\begin{figure}[htbp]
\centering
\includegraphics[width=0.55\textwidth]{figures/3ph_semi_ckt.png}
\caption{Three-Phase Semi-Converter: Bridge topology with 3 thyristors in positive group and 3 diodes in negative group.}
\end{figure}

\begin{figure}[htbp]
\centering
\includegraphics[width=0.48\textwidth]{figures/3ph_semi_wave_a0.png}
\hfill
\includegraphics[width=0.48\textwidth]{figures/3ph_semi_wave_a60.png}
\caption{Three-Phase Semi-Converter Waveforms: (Left) $\alpha = 0^\circ$ (6-pulse operation). (Right) $\alpha = 60^\circ$ (transition to 3-pulse freewheeling operation).}
\end{figure}

\subsubsection{The Critical Transition at \texorpdfstring{$\alpha = 60^\circ$}{alpha = 60 deg}}
The operation of the 3-phase semi-converter bifurcates into two completely distinct operational regimes:
\begin{enumerate}[leftmargin=*]
  \item \textbf{Continuous 6-Pulse Conduction ($\alpha \le 60^\circ$):}
  The output voltage never touches zero. At any instant, one thyristor from the top group and one diode from the bottom group (belonging to different phases) conduct. The output voltage is composed of 6 distinct line voltage segments per cycle, with a ripple frequency of $6 f_s = 300\text{ Hz}$. Freewheeling does not occur because the instantaneous output line voltage remains strictly positive.
  \item \textbf{Freewheeling 3-Pulse Conduction ($\alpha > 60^\circ$):}
  As $\alpha$ exceeds $60^\circ$, the instantaneous line voltage tries to swing negative. The moment it reaches zero, the diode in the same phase leg as the conducting thyristor turns ON. Current now circulates (freewheels) internally through that thyristor and diode pair in the same leg! 
  
  During this freewheeling interval, the load terminals are short-circuited through the leg, clamping the output voltage to \textbf{zero} ($v_o = 0$) and disconnecting the AC line current ($i_s = 0$). Three line voltage pulses disappear from the output envelope, reducing the ripple frequency to $3 f_s = 150\text{ Hz}$.
\end{enumerate}

\subsubsection{Derivation of Average Output Voltage (\texorpdfstring{$V_o$}{Vo})}
Each of the 3 thyristors conducts for $120^\circ$ ($2\pi/3$). Taking the line voltage as reference ($v_{line} = V_{ml}\sin\omega t$), active conduction delivers voltage from $\omega t = \alpha$ up to $\pi$:
\[
V_o = \frac{3}{2\pi} \int_{\alpha}^{\pi} V_{ml} \sin\omega t \, d(\omega t) = \frac{3 V_{ml}}{2\pi} [-\cos\omega t]_\alpha^\pi = \frac{3 V_{ml}}{2\pi} [-\cos\pi - (-\cos\alpha)]
\]
\begin{formulabox}[3-Phase Semi-Converter Output Voltage]
\[
\boxed{V_o = \frac{3 V_{ml}}{2\pi} (1 + \cos\alpha) = \frac{3\sqrt{3} V_{mp}}{2\pi} (1 + \cos\alpha) = \frac{3\sqrt{2} V_L}{2\pi} (1 + \cos\alpha) \approx 0.675 V_L (1 + \cos\alpha)}
\]
\begin{itemize}[leftmargin=*]
  \item At $\alpha = 0^\circ$: $V_o = \frac{3 V_{ml}}{\pi} \approx 1.35 V_L$ (identical to 6-pulse full converter).
  \item At $\alpha = 60^\circ$: $V_o = \frac{3 V_{ml}}{2\pi}(1 + 0.5) = \frac{4.5 V_{ml}}{2\pi} \approx 1.013 V_L$.
  \item At $\alpha = 180^\circ$: $V_o = 0\text{ V}$.
\end{itemize}
\end{formulabox}

\subsection{Source Inductance in Three-Phase Full Converter}

\textbf{Physical Commutation Process in 3-Phase Bridges:}
In a 3-phase full converter, commutation occurs \textbf{6 times per AC supply cycle} (once every $60^\circ$). When an incoming thyristor (say $T_3$ in phase $b$) is fired at $\alpha$, the outgoing thyristor ($T_1$ in phase $a$) cannot shut off instantaneously because the supply inductance $L_s$ in phase $a$ opposes current collapse ($v_{Ls} = L_s di/dt$).

\begin{conceptbox}[Commutation Overlap Mechanics]
During the overlap interval $\mu$:
\begin{itemize}[leftmargin=*]
  \item \textbf{Three devices conduct simultaneously:} The two commutating devices in one group ($T_1$ and $T_3$) plus the one return device in the opposite group ($T_2$ in phase $c$).
  \item \textbf{Phase Line Short-Circuit:} Phases $a$ and $b$ are tied together through their source inductances $L_s$. The potential of the positive DC rail is the midpoint voltage: $v_{pos} = \frac{v_a + v_b}{2}$.
  \item The output line voltage during overlap is reduced by $\frac{v_a - v_b}{2} = \frac{v_{ab}}{2}$.
  \item This loss of volt-second area causes a net reduction $\Delta V_o$ in the average DC output voltage.
\end{itemize}
\end{conceptbox}

\subsubsection{Derivation of Overlap Angle \texorpdfstring{$\mu$}{mu} and Voltage Drop \texorpdfstring{$\Delta V_o$}{Delta Vo}}
During overlap of phase $a$ and phase $b$:
\[
v_a - L_s \frac{di_a}{dt} = v_b - L_s \frac{di_b}{dt} \implies v_{ab} = L_s \left( \frac{di_a}{dt} - \frac{di_b}{dt} \right)
\]
Since $i_a + i_b = I_o$ (constant DC current), $\frac{di_a}{dt} = -\frac{di_b}{dt}$. Thus:
\[
v_{ab} = 2 L_s \frac{di_b}{dt} = V_{ml} \sin\omega t
\]
Integrating incoming current $i_b$ from $0$ to $I_o$ as $\omega t$ goes from $\alpha$ to $\alpha+\mu$:
\[
\int_{0}^{I_o} di_b = \frac{V_{ml}}{2\omega L_s} \int_{\alpha}^{\alpha+\mu} \sin\omega t \, d(\omega t) \implies I_o = \frac{V_{ml}}{2\omega L_s} [\cos\alpha - \cos(\alpha+\mu)]
\]
\begin{formulabox}[3-Phase Overlap Angle \& Equivalent Resistance]
\[
\boxed{\cos\alpha - \cos(\alpha+\mu) = \frac{2\omega L_s I_o}{V_{ml}} = \frac{\sqrt{2}\omega L_s I_o}{V_L}}
\]
Because there are \textbf{6 commutations per cycle}, the average voltage drop across the period $2\pi$ is:
\[
\Delta V_o = \frac{6 \times (\text{lost volt-second area per commutation})}{2\pi} = \frac{6 \times (\omega L_s I_o)}{2\pi} = \boxed{\frac{3\omega L_s}{\pi} I_o = 3 f L_s I_o}
\]
Notice that in 3-phase, $\Delta V_o = 3 f L_s I_o$, whereas in 1-phase it is $2 f L_s I_o$ (proportional to the number of commutations per cycle!).
\[
\boxed{R_c = \frac{3\omega L_s}{\pi} = 3 f L_s \quad (\Omega)}
\]
\[
\boxed{V_o = \frac{3 V_{ml}}{\pi}\cos\alpha - \Delta V_o = \frac{3\sqrt{2} V_L}{\pi}\cos\alpha - 3 f L_s I_o = \frac{3 V_{ml}}{2\pi}[\cos\alpha + \cos(\alpha+\mu)]}
\]
\end{formulabox}

%%%%%%%%%%%%%%%%%%%%%%%%%%%%%%%%%%%%%%%%%%%%%%%%%%%%%%%%%%%%%%%%%%%%%%%%%%%%%%%
\section{DC-DC Switching Regulators (Buck, Boost, Buck-Boost)}
%%%%%%%%%%%%%%%%%%%%%%%%%%%%%%%%%%%%%%%%%%%%%%%%%%%%%%%%%%%%%%%%%%%%%%%%%%%%%%%

Switch-mode converters convert an unregulated DC input $V_d$ into a regulated output $V_o$ with theoretical efficiency of $100\%$ by using switches in saturation or cutoff.

\begin{figure}[htbp]
\centering
\includegraphics[width=0.48\textwidth]{figures/dcdc_block.png}
\hfill
\includegraphics[width=0.48\textwidth]{figures/dcdc_pwm.png}
\caption{DC-DC Switch-Mode Conversion: (Left) Functional block diagram. (Right) Pulse-Width Modulation (PWM) duty cycle control by comparing control voltage with sawtooth carrier.}
\end{figure}

\begin{figure}[htbp]
\centering
\includegraphics[width=0.48\textwidth]{figures/dcdc_linear.png}
\hfill
\includegraphics[width=0.48\textwidth]{figures/dcdc_smps.png}
\caption{Power Supply Comparison: (Left) Linear regulator dissipates $(V_{\text{in}} - V_o)I_o$ as continuous heat. (Right) Switching power supply operates transistor as an efficient high-frequency switch.}
\end{figure}

\subsection{Two Fundamental Steady-State Laws}
\begin{formulabox}[Conservation Principles in Power Electronics]
\begin{enumerate}[leftmargin=*]
  \item \textbf{Inductor Volt-Second Balance:} Because $v_L = L \frac{di_L}{dt}$ and $i_L(T_s) = i_L(0)$:
  \[
  \int_{0}^{T_s} v_L(t) \, dt = 0 \iff \langle v_L \rangle = 0
  \]
  \item \textbf{Capacitor Charge Balance:} Because $i_C = C \frac{dv_C}{dt}$ and $v_C(T_s) = v_C(0)$:
  \[
  \int_{0}^{T_s} i_C(t) \, dt = 0 \iff \langle i_C \rangle = 0
  \]
\end{enumerate}
\end{formulabox}

\begin{conceptbox}[The Small-Ripple Approximation (Key Analytical Assumption)]
In steady-state CCM operation, the inductor current $i_L(t)$ and capacitor voltage $v_C(t)$ contain small triangular AC ripples superimposed on large DC values. The \textbf{small-ripple approximation} assumes:
\[
i_L(t) \approx I_L \quad (\text{average value}), \qquad v_C(t) \approx V_o \quad (\text{average value})
\]
This allows us to treat $i_L$ and $v_C$ as constants when writing KVL/KCL during the ON and OFF intervals, greatly simplifying the analysis. This approximation is valid when $\Delta I_L \ll I_L$ and $\Delta V_o \ll V_o$, which holds whenever the filter $L$ and $C$ are properly designed.
\end{conceptbox}

\subsection{Step-Down (Buck) Converter}
\begin{figure}[htbp]
\centering
\includegraphics[width=0.55\textwidth]{figures/dcdc_buck_ckt.png}
\caption{Step-Down (Buck) Converter: Basic circuit topology consisting of power switch $S$, freewheeling diode $D$, and output low-pass filter $L-C$.}
\end{figure}

\begin{figure}[htbp]
\centering
\includegraphics[width=0.58\textwidth]{figures/dcdc_buck_wave.png}
\caption{Buck Converter CCM Waveforms: Switch state, diode voltage $v_D$, inductor voltage $v_L$, triangular inductor current $i_L$, capacitor current $i_C$, and AC ripple voltage $\Delta v_o$.}
\end{figure}

\begin{figure}[htbp]
\centering
\includegraphics[width=0.52\textwidth]{figures/dcdc_buck_ripple.png}
\caption{Buck Converter Output Voltage Ripple: Geometric charge integration of the positive shaded triangular area $\Delta Q = \frac{T_s \Delta I_L}{8}$.}
\end{figure}

\subsubsection{Circuit Operation \& Sub-Interval Analysis}
The Buck converter switches between two circuit states at high frequency $f_s = 1/T_s$:
\begin{itemize}[leftmargin=*]
  \item \textbf{State 1: Switch ON ($0 \le t \le D T_s$):}
  Switch $S$ is closed; diode $D$ is reverse-biased by input voltage $V_d$. 
  The inductor is connected between input $V_d$ and output $V_o$. Applying KVL:
  \[
  v_L(t) = V_d - V_o > 0
  \]
  Because $v_L > 0$, the inductor current ramps \textbf{upward} linearly with constant positive slope:
  \[
  \frac{di_L}{dt} = \frac{V_d - V_o}{L}
  \]
  During this time, the input DC source simultaneously feeds the load and stores magnetic energy in the inductor.
  
  \item \textbf{State 2: Switch OFF ($D T_s \le t \le T_s$):}
  Switch $S$ is opened. The collapsing magnetic field of the inductor forces current to continue flowing in the same direction, which forward-biases freewheeling diode $D$. 
  The diode clamps the left side of the inductor to $0\text{ V}$ (ground). Applying KVL:
  \[
  v_L(t) = 0 - V_o = -V_o < 0
  \]
  Because $v_L < 0$, the inductor current ramps \textbf{downward} linearly with constant negative slope:
  \[
  \frac{di_L}{dt} = -\frac{V_o}{L}
  \]
  During this interval, the inductor discharges its stored magnetic energy into the load capacitor circuit.
\end{itemize}

\subsubsection{Derivation of Voltage Gain (Volt-Second Balance)}
In steady state, the average voltage across an ideal inductor over one complete switching period $T_s$ must be exactly zero:
\[
\int_{0}^{T_s} v_L(t) \, dt = \underbrace{(V_d - V_o)(D T_s)}_{\text{Volt-seconds during ON}} + \underbrace{(-V_o)((1 - D) T_s)}_{\text{Volt-seconds during OFF}} = 0
\]
Dividing through by $T_s$:
\[
(V_d - V_o)D - V_o(1 - D) = 0 \implies V_d D - V_o D - V_o + V_o D = 0
\]
\begin{formulabox}[Buck Converter Voltage Conversion Ratio]
\[
\boxed{V_o = D V_d}
\]
Since the duty cycle is constrained to $0 < D < 1$, the output voltage is always strictly less than the input voltage ($V_o < V_d$). Hence the name \textbf{Step-Down (Buck)}.
\end{formulabox}

\subsubsection{Inductor Ripple \& Critical Inductance \texorpdfstring{$L_{\min}$}{Lmin} for CCM}
The peak-to-peak inductor current ripple $\Delta I_L$ is the net rise in current during the ON interval:
\[
\Delta I_L = \left(\frac{di_L}{dt}\right) \times (D T_s) = \frac{V_d - V_o}{L} D T_s = \frac{V_d(1 - D)}{L} D T_s = \boxed{\frac{V_o (1 - D)}{f_s L}}
\]
The average inductor current equals the DC load current: $I_L = I_o = \frac{V_o}{R}$. 

For \textbf{Continuous Conduction Mode (CCM)}, the inductor current must never touch zero. This requires the valley current to be positive:
\[
I_{L,\min} = I_L - \frac{\Delta I_L}{2} \ge 0 \implies I_L \ge \frac{\Delta I_L}{2}
\]
At the boundary between CCM and Discontinuous Conduction Mode (DCM):
\[
\frac{V_o}{R} = \frac{1}{2} \left[ \frac{V_o (1 - D)}{f_s L_{\min}} \right] \implies \boxed{L_{\min} = \frac{(1 - D) R}{2 f_s}}
\]
If $L < L_{\min}$, the inductor current drops to zero before the end of the period (DCM), causing the output voltage to rise above $D V_d$ and become load-dependent.

\subsubsection{Output Voltage Ripple \& Capacitor Sizing \texorpdfstring{$C_{\min}$}{Cmin}}
Refer to Figure 24 showing the capacitor current $i_C(t) = i_L(t) - I_o$. The capacitor absorbs current whenever $i_L(t) > I_o$ (the upper shaded triangle of height $\Delta I_L / 2$ and base width $T_s / 2$). The total charge accumulated during this half-cycle is the area of that triangle:
\[
\Delta Q = \frac{1}{2} \times \left(\frac{T_s}{2}\right) \times \left(\frac{\Delta I_L}{2}\right) = \frac{T_s \Delta I_L}{8} = \frac{\Delta I_L}{8 f_s}
\]
Since $\Delta V_o = \frac{\Delta Q}{C}$:
\[
\Delta V_o = \frac{\Delta I_L}{8 C f_s} = \frac{\frac{V_o(1 - D)}{f_s L}}{8 C f_s} = \frac{V_o (1 - D)}{8 L C f_s^2}
\]
\begin{formulabox}[Buck Capacitor Sizing Equation]
\[
\boxed{\frac{\Delta V_o}{V_o} = \frac{1 - D}{8 L C f_s^2} \implies C_{\min} = \frac{1 - D}{8 L f_s^2 (\Delta V_o / V_o)}}
\]
Notice the $f_s^2$ term in the denominator: doubling the switching frequency reduces the required filter capacitance by a factor of 4!
\end{formulabox}

\subsection{Step-Up (Boost) Converter}
\begin{figure}[htbp]
\centering
\includegraphics[width=0.55\textwidth]{figures/dcdc_boost_ckt.png}
\caption{Step-Up (Boost) Converter: Circuit topology showing inductor in series with the input DC source.}
\end{figure}

\begin{figure}[htbp]
\centering
\includegraphics[width=0.58\textwidth]{figures/dcdc_boost_wave.png}
\caption{Boost Converter CCM Waveforms: Switch control signal, inductor voltage $v_L$, continuous inductor current $i_L$, discontinuous diode current $i_D$, and capacitor current $i_C$.}
\end{figure}

\subsubsection{How Does the Voltage Get Stepped Up?}
A beginner often wonders: \textit{``How can a passive inductor output a higher DC voltage than the input battery?''}
The key lies in the two-stage energy transfer:
\begin{itemize}[leftmargin=*]
  \item \textbf{Stage 1: Charging the Inductor ($0 \le t \le D T_s$):}
  Switch $S$ is closed. The diode $D$ is reverse-biased because cathode is at $V_o > 0$. The inductor is tied directly across input voltage $V_d$:
  \[
  v_L(t) = V_d > 0 \implies \frac{di_L}{dt} = \frac{V_d}{L}
  \]
  Energy is drawn from the input source and stored in the inductor's magnetic field. Meanwhile, the output capacitor $C$ alone supplies the entire load current $I_o$.
  
  \item \textbf{Stage 2: Boosting \& Discharging ($D T_s \le t \le T_s$):}
  Switch $S$ opens. Inductor current cannot drop instantly! The collapsing magnetic field causes the inductor terminal voltage to flip polarity (left side positive, right side higher positive). 
  The inductor now acts as an \textbf{auxiliary booster battery connected in series with the input source}:
  \[
  v_L(t) = V_d - V_o < 0
  \]
  Because the switch terminal reaches $V_o$, the diode conducts, dumping current from \textit{both} the input source and the discharging inductor simultaneously into the output capacitor and load. The resulting output voltage is necessarily higher than the input: $V_o > V_d$!
\end{itemize}

\subsubsection{Volt-Second Balance \& Derivations}
Applying Inductor Volt-Second Balance:
\[
\int_{0}^{T_s} v_L(t) \, dt = (V_d)(D T_s) + (V_d - V_o)((1 - D) T_s) = 0
\]
\[
V_d D + (V_d - V_o)(1 - D) = 0 \implies V_d [D + 1 - D] = V_o (1 - D)
\]
\begin{formulabox}[Boost Converter Voltage Gain]
\[
\boxed{V_o = \frac{V_d}{1 - D}}
\]
Because $0 < D < 1$, the factor $\frac{1}{1 - D} > 1$, meaning $V_o > V_d$ for all duty cycles.
\end{formulabox}

\subsubsection{Inductor Current \& Sizing Equations}
Assuming an ideal converter (100\% efficiency), input power equals output power:
\[
P_{\text{in}} = P_{\text{out}} \implies V_d I_{\text{in}} = V_o I_o
\]
Since the input source is in direct series with the inductor, the average input current is the average inductor current: $I_{\text{in}} = I_L$. Therefore:
\[
I_L = \frac{V_o}{V_d} I_o = \frac{1}{1 - D} I_o = \boxed{\frac{V_d}{(1 - D)^2 R}}
\]
\begin{itemize}[leftmargin=*]
  \item Inductor Current Ripple: $\Delta I_L = \frac{V_d D T_s}{L} = \boxed{\frac{V_d D}{f_s L}}$
  \item Boundary Condition for CCM ($I_L \ge \Delta I_L / 2$):
  \[
  \frac{V_d}{(1 - D)^2 R} \ge \frac{1}{2} \left[ \frac{V_d D}{f_s L_{\min}} \right] \implies \boxed{L_{\min} = \frac{D (1 - D)^2 R}{2 f_s}}
  \]
  \item Capacitor Ripple \& Sizing: During the switch ON interval ($D T_s$), the diode is OFF, so the capacitor must supply the full load current $I_o = V_o/R$. The total charge removed from $C$ is $\Delta Q = I_o (D T_s)$. Therefore:
  \[
  \Delta V_o = \frac{\Delta Q}{C} = \frac{I_o D T_s}{C} = \frac{V_o D}{R C f_s} \implies \boxed{\frac{\Delta V_o}{V_o} = \frac{D}{R C f_s} \implies C_{\min} = \frac{D}{R f_s (\Delta V_o / V_o)}}
  \]
\end{itemize}

\subsection{Step-Down/Step-Up (Buck-Boost) Converter}
\begin{figure}[htbp]
\centering
\includegraphics[width=0.52\textwidth]{figures/dcdc_buckboost_ckt.png}
\caption{Step-Down/Step-Up (Buck-Boost) Converter: Circuit schematic showing reversed diode polarity yielding negative output voltage.}
\end{figure}

\begin{figure}[htbp]
\centering
\includegraphics[width=0.58\textwidth]{figures/dcdc_buckboost_wave.png}
\caption{Buck-Boost Converter CCM Waveforms: Switch signal, inductor voltage $v_L$, inductor current $i_L$, diode current, and capacitor current.}
\end{figure}

\subsubsection{Circuit Operation \& Why Is Output Polarity Inverted?}
Look closely at the schematic in Figure 27: the diode cathode is connected to the inductor, and its anode is connected to the output DC rail!
\begin{itemize}[leftmargin=*]
  \item \textbf{Switch ON ($0 \le t \le D T_s$):}
  Switch $S$ is closed. The diode $D$ is reverse-biased. The inductor connects directly across the input source $V_d$, drawing current from top to bottom:
  \[
  v_L(t) = V_d > 0 \implies \frac{di_L}{dt} = \frac{V_d}{L}
  \]
  Energy is stored entirely in the inductor. The load is completely isolated from the source and is powered solely by capacitor $C$.
  
  \item \textbf{Switch OFF ($D T_s \le t \le T_s$):}
  Switch $S$ opens. The inductor current \textbf{must continue flowing in the same downward direction} (into ground). To sustain this downward flow, current is pulled out of the bottom ground rail and pushed \textit{upward} into the anode of diode $D$!
  
  This forces the top rail of the load to become \textbf{negative with respect to ground}. Hence, the output voltage is inverted:
  \[
  v_L(t) = V_o < 0 \quad (\text{negative relative to ground})
  \]
\end{itemize}

\subsubsection{Volt-Second Balance \& Sizing Formulae}
Applying Inductor Volt-Second Balance:
\[
(V_d)(D T_s) + (V_o)((1 - D) T_s) = 0 \implies \boxed{V_o = -\frac{D}{1 - D} V_d \implies |V_o| = \frac{D}{1 - D} V_d}
\]
\begin{formulabox}[Buck-Boost Operational Regimes]
\begin{itemize}[leftmargin=*]
  \item For $D < 0.5 \implies \frac{D}{1 - D} < 1 \implies |V_o| < V_d$ (\textbf{Step-Down / Buck Mode}).
  \item For $D = 0.5 \implies \frac{D}{1 - D} = 1 \implies |V_o| = V_d$ (Unity gain inverter).
  \item For $D > 0.5 \implies \frac{D}{1 - D} > 1 \implies |V_o| > V_d$ (\textbf{Step-Up / Boost Mode}).
\end{itemize}
\end{formulabox}

\subsubsection{Component Sizing Equations}
\begin{itemize}[leftmargin=*]
  \item Average Inductor Current: $I_L = I_{\text{in}} + I_o = \frac{D}{1 - D} I_o + I_o = \boxed{\frac{I_o}{1 - D} = \frac{D V_d}{(1 - D)^2 R}}$
  \item Inductor Current Ripple: $\Delta I_L = \frac{V_d D}{f_s L}$
  \item Critical Inductance for CCM ($I_L \ge \Delta I_L / 2$):
  \[
  \frac{|V_o|}{(1 - D) R} \ge \frac{1}{2} \left[ \frac{V_d D}{f_s L_{\min}} \right] \implies \boxed{L_{\min} = \frac{(1 - D)^2 R}{2 f_s}}
  \]
  \item Output Voltage Ripple: During switch ON, diode is OFF and capacitor supplies the full load current $I_o$:
  \[
  \Delta Q = I_o (D T_s) \implies \boxed{\frac{\Delta V_o}{|V_o|} = \frac{D}{R C f_s} \implies C_{\min} = \frac{D}{R f_s (\Delta V_o / |V_o|)}}
  \]
\end{itemize}

\subsection{Effect of Non-Idealities}
\begin{alertbox}[Inductor Resistance \texorpdfstring{$r_L$}{rL} in Boost Converter]
Due to inductor winding resistance $r_L$, the actual voltage gain of a Boost converter is:
\[
\boxed{\frac{V_o}{V_d} = \frac{1}{1 - D} \cdot \frac{1}{1 + \frac{r_L}{R (1 - D)^2}}}
\]
As $D \to 1$, the inductor current $I_L = I_o / (1-D)$ becomes massive, so $I_L^2 r_L$ losses completely dominate! The voltage gain peaks around $D \approx 0.85$ and drops sharply to zero at $D = 1$.
\end{alertbox}

\begin{alertbox}[Capacitor ESR (\texorpdfstring{$r_C$}{rC})]
Capacitor Equivalent Series Resistance produces an additional voltage step: $\Delta V_{\text{ESR}} = \Delta I \cdot r_C$. At switching frequencies above $50\text{ kHz}$, ESR ripple often exceeds pure capacitive ripple!
\end{alertbox}

\subsection{Selection of Power Semiconductor Devices}
\begin{conceptbox}[Device Selection Guidelines for DC-DC Converters]
\begin{itemize}[leftmargin=*]
  \item \textbf{Power MOSFET:} Preferred for low-to-medium voltage ($< 500\text{ V}$), high-frequency ($50\text{ kHz}$ to several MHz) applications. Very fast switching transitions ($< 100\text{ ns}$) minimise switching losses. Conduction loss is $I_D^2 \cdot R_{DS(on)}$, which increases with voltage rating. Ideal choice for \textbf{Buck converters} at $V_d < 200\text{ V}$.
  \item \textbf{IGBT (Insulated Gate Bipolar Transistor):} Preferred for medium-to-high voltage ($> 600\text{ V}$), moderate-frequency ($< 50\text{ kHz}$) applications. Lower conduction loss than MOSFET at high voltages due to bipolar conductivity modulation. Limited by \textit{current tailing} during turn-off. Ideal for \textbf{Boost/Buck-Boost converters} at higher power levels.
  \item \textbf{Schottky Diode:} Negligible reverse recovery time ($t_{rr} \approx 0$) and low forward voltage drop ($V_F \approx 0.3\text{--}0.5\text{ V}$). Ideal freewheeling diode in low-voltage SMPS ($V_o < 100\text{ V}$) where recovery losses of $p$-$n$ junction diodes would be unacceptable.
  \item \textbf{Fast Recovery Diode (FRD):} Used as freewheeling diode in higher voltage ($> 200\text{ V}$) converters where Schottky diodes are unavailable. Has short but non-zero $t_{rr}$ ($< 100\text{ ns}$).
\end{itemize}
\textbf{Key Trade-off:} Higher switching frequency $f_s$ reduces filter $L$ and $C$ sizes (smaller, lighter converter), but increases switching losses. The device choice must balance these factors.
\end{conceptbox}

\newpage
%%%%%%%%%%%%%%%%%%%%%%%%%%%%%%%%%%%%%%%%%%%%%%%%%%%%%%%%%%%%%%%%%%%%%%%%%%%%%%%
\section{Comprehensive Solved Numerical Problems}
%%%%%%%%%%%%%%%%%%%%%%%%%%%%%%%%%%%%%%%%%%%%%%%%%%%%%%%%%%%%%%%%%%%%%%%%%%%%%%%

\begin{examplebox}[Problem 1: 1-Phase Full Converter RL Load Performance]
\textbf{Question:} A single-phase full converter operates from $V_s = 230\text{ V}, 50\text{ Hz}$ supply with $\alpha = 45^\circ$ feeding a ripple-free current $I_o = 15\text{ A}$. Calculate $V_o, V_{o,\text{rms}}, P, I_{s1}, DF, g, PF,$ and $THD$.

\tcbline
\textbf{Solution:}
\begin{enumerate}[leftmargin=*]
  \item Peak voltage: $V_m = \sqrt{2} \times 230 = 325.27\text{ V}$.
  \item Average Output Voltage:
  \[
  V_o = \frac{2 V_m}{\pi} \cos\alpha = \frac{2 \times 325.27}{\pi} \cos(45^\circ) = 207.07 \times 0.7071 = \mathbf{146.42\text{ V}}
  \]
  \item RMS Output Voltage: $V_{o,\text{rms}} = V_s = \mathbf{230\text{ V}}$.
  \item Real Power: $P = V_o I_o = 146.42 \times 15 = \mathbf{2196.3\text{ W}}$.
  \item Fundamental RMS Current: $I_{s1} = \frac{2\sqrt{2}}{\pi} I_o = 0.9003 \times 15 = \mathbf{13.505\text{ A}}$.
  \item Displacement Factor: $DF = \cos(45^\circ) = \mathbf{0.7071\text{ (lagging)}}$.
  \item Distortion Factor: $g = \frac{I_{s1}}{I_s} = \frac{13.505}{15} = \mathbf{0.9003}$.
  \item Input Power Factor: $PF = g \times DF = 0.9003 \times 0.7071 = \mathbf{0.6366\text{ (lagging)}}$.
  \item Current THD: $THD = \sqrt{\frac{1}{g^2}-1} = \sqrt{\frac{\pi^2}{8}-1} = \mathbf{48.34\%}$.
\end{enumerate}
\end{examplebox}

\begin{examplebox}[Problem 2: 1-Phase Full Converter Inverter Mode]
\textbf{Question:} A $230\text{ V}, 50\text{ Hz}$ full converter operates in inverter mode to regenerate power into the AC grid from a $150\text{ V}$ battery ($R = 1.5\ \Omega$) with $I_o = 12\text{ A}$. Find $\alpha$ and power regenerated.

\tcbline
\textbf{Solution:}
\begin{enumerate}[leftmargin=*]
  \item Loop KVL in inverter mode:
  \[
  V_o = -(E - I_o R) = -(150 - 12 \times 1.5) = -(150 - 18) = \mathbf{-132\text{ V}}
  \]
  \item Since $V_o = \frac{2 V_m}{\pi} \cos\alpha$:
  \[
  -132 = 207.07 \cos\alpha \implies \cos\alpha = \frac{-132}{207.07} = -0.63746
  \]
  \[
  \alpha = \cos^{-1}(-0.63746) = \mathbf{129.6^\circ}
  \]
  \item Power delivered to AC source: $P_{\text{regen}} = |V_o| I_o = 132 \times 12 = \mathbf{1584\text{ W}} = \mathbf{1.584\text{ kW}}$.
\end{enumerate}
\end{examplebox}

\begin{examplebox}[Problem 3: Semi-Converter vs. Full Converter Power Factor]
\textbf{Question:} A single-phase semi-converter operates from $230\text{ V}, 50\text{ Hz}$ at $\alpha = 60^\circ$ with $I_o = 10\text{ A}$. Calculate $V_o, I_s, DF,$ and $PF$. Compare with a full converter.

\tcbline
\textbf{Solution:}
\begin{enumerate}[leftmargin=*]
  \item $V_o = \frac{V_m}{\pi}(1 + \cos 60^\circ) = \frac{325.27}{\pi} (1.5) = \mathbf{155.31\text{ V}}$.
  \item $I_s = I_o \sqrt{\frac{180^\circ - 60^\circ}{180^\circ}} = 10 \times \sqrt{2/3} = \mathbf{8.165\text{ A}}$.
  \item $DF = \cos(60^\circ / 2) = \cos(30^\circ) = \mathbf{0.8660}$.
  \item $P = V_o I_o = 1553.1\text{ W}; \quad S = 230 \times 8.165 = 1877.95\text{ VA} \implies PF = \frac{1553.1}{1877.95} = \mathbf{0.8270}$.
  \item Full converter at $60^\circ$: $PF_{\text{full}} = 0.9003 \cos(60^\circ) = \mathbf{0.4502}$.
  \[
  \frac{PF_{\text{semi}}}{PF_{\text{full}}} = \frac{0.8270}{0.4502} \approx \mathbf{1.84} \quad (\mathbf{84\% \text{ higher!}})
  \]
\end{enumerate}
\end{examplebox}

\begin{examplebox}[Problem 4: Source Inductance in 1-Phase Full Converter]
\textbf{Question:} $V_s = 230\text{ V}, 50\text{ Hz}, L_s = 2.5\text{ mH}, I_o = 20\text{ A}, \alpha = 30^\circ$. Find $\mu, \Delta V_o, V_o,$ and $R_c$.

\tcbline
\textbf{Solution:}
\begin{enumerate}[leftmargin=*]
  \item $\frac{2\omega L_s I_o}{V_m} = \frac{2 \times 314.16 \times 0.0025 \times 20}{325.27} = 0.09658$.
  \item $\cos(30^\circ + \mu) = \cos(30^\circ) - 0.09658 = 0.86603 - 0.09658 = 0.76945$.
  \[
  30^\circ + \mu = \cos^{-1}(0.76945) = 39.70^\circ \implies \mu = \mathbf{9.70^\circ}
  \]
  \item $\Delta V_o = 2 f L_s I_o = 2 \times 50 \times 0.0025 \times 20 = \mathbf{5.0\text{ V}}$.
  \item $V_{o,\text{ideal}} = 207.07 \cos(30^\circ) = 179.33\text{ V} \implies V_o = 179.33 - 5.0 = \mathbf{174.33\text{ V}}$.
  \item $R_c = 2 f L_s = \mathbf{0.25\ \Omega}$.
\end{enumerate}
\end{examplebox}

\begin{examplebox}[Problem 5: 3-Phase Full Converter Feeding SEDC Motor]
\textbf{Question:} A 3-phase full converter operates from $400\text{ V}$ line, $50\text{ Hz}$. It feeds a DC motor with $E = 350\text{ V}, R = 0.5\ \Omega, I_o = 40\text{ A}$. Find $\alpha, P,$ and input $PF$.

\tcbline
\textbf{Solution:}
\begin{enumerate}[leftmargin=*]
  \item Required terminal voltage: $V_o = E + I_o R = 350 + (40 \times 0.5) = \mathbf{370\text{ V}}$.
  \item Peak line voltage: $V_{ml} = \sqrt{2} \times 400 = 565.69\text{ V}$.
  \[
  V_o = \frac{3 V_{ml}}{\pi} \cos\alpha \implies 370 = \frac{3 \times 565.69}{\pi} \cos\alpha = 540.19 \cos\alpha
  \]
  \[
  \cos\alpha = \frac{370}{540.19} = 0.6849 \implies \alpha = \mathbf{46.77^\circ}
  \]
  \item Active Power: $P = V_o I_o = 370 \times 40 = \mathbf{14.8\text{ kW}}$.
  \item Input Power Factor: $PF = \frac{3}{\pi} \cos\alpha = 0.9549 \times 0.6849 = \mathbf{0.654\text{ (lagging)}}$.
\end{enumerate}
\end{examplebox}

\begin{examplebox}[Problem 6: 3-Phase Semi-Converter]
\textbf{Question:} Repeat Problem 5 if a 3-phase semi-converter is used.

\tcbline
\textbf{Solution:}
\[
V_o = \frac{3 V_{ml}}{2\pi} (1 + \cos\alpha) \implies 370 = 270.095 (1 + \cos\alpha)
\]
\[
1 + \cos\alpha = \frac{370}{270.095} = 1.36989 \implies \cos\alpha = 0.36989 \implies \alpha = \mathbf{68.29^\circ}
\]
Since $\alpha > 60^\circ$, it operates in the freewheeling 3-pulse regime.
\end{examplebox}

\begin{examplebox}[Problem 7: 3-Phase Full Converter with Source Inductance]
\textbf{Question:} $V_L = 415\text{ V}, 50\text{ Hz}, L_s = 1.2\text{ mH}, I_o = 50\text{ A}, \alpha = 45^\circ$. Find $\mu, \Delta V_o,$ and $V_o$.

\tcbline
\textbf{Solution:}
\begin{enumerate}[leftmargin=*]
  \item $V_{ml} = \sqrt{2} \times 415 = 586.9\text{ V}$.
  \[
  \cos\alpha - \cos(\alpha+\mu) = \frac{2\omega L_s I_o}{\sqrt{3} V_{ml}} = \frac{2 \times 314.16 \times 0.0012 \times 50}{\sqrt{3} \times 586.9} = 0.037085
  \]
  \[
  \cos(45^\circ + \mu) = \cos(45^\circ) - 0.037085 = 0.70711 - 0.037085 = 0.67002
  \]
  \[
  45^\circ + \mu = 47.93^\circ \implies \mu = \mathbf{2.93^\circ}
  \]
  \item Voltage Drop: $\Delta V_o = 3 f L_s I_o = 3 \times 50 \times 0.0012 \times 50 = \mathbf{9.0\text{ V}}$.
  \item Ideal: $V_{o,\text{ideal}} = 560.46 \cos(45^\circ) = 396.31\text{ V} \implies V_o = 396.31 - 9.0 = \mathbf{387.31\text{ V}}$.
\end{enumerate}
\end{examplebox}

\begin{examplebox}[Problem 8: Complete Buck Converter Design]
\textbf{Question:} $V_d = 48\text{ V}, V_o = 12\text{ V}, R = 6\ \Omega, f_s = 40\text{ kHz}$. Find $D, L_{\min}$, actual ripples for $L = 1.5 L_{\min}$, and $C$ for $\Delta V_o / V_o \le 0.5\%$.

\tcbline
\textbf{Solution:}
\begin{enumerate}[leftmargin=*]
  \item $D = \frac{V_o}{V_d} = \frac{12}{48} = \mathbf{0.25}$.
  \item $L_{\min} = \frac{(1 - D) R}{2 f_s} = \frac{(1 - 0.25) \times 6}{2 \times 40,000} = \mathbf{56.25\ \mu\text{H}}$.
  \item $L = 1.5 \times 56.25 = 84.375\ \mu\text{H}$.
  \[
  \Delta I_L = \frac{V_o (1 - D)}{f_s L} = \frac{12 \times 0.75}{40,000 \times 84.375 \times 10^{-6}} = \mathbf{2.667\text{ A}}
  \]
  \[
  I_{L,\max} = I_o + \frac{\Delta I_L}{2} = 2 + 1.333 = \mathbf{3.333\text{ A}}
  \]
  \item Filter Capacitor:
  \[
  C_{\min} = \frac{1 - D}{8 L f_s^2 (\Delta V_o / V_o)} = \frac{0.75}{8 \times 84.375 \times 10^{-6} \times (40,000)^2 \times 0.005} = \mathbf{138.9\ \mu\text{F}}
  \]
\end{enumerate}
\end{examplebox}

\begin{examplebox}[Problem 9: Boost Converter with Inductor Resistance \texorpdfstring{$r_L$}{rL}]
\textbf{Question:} $V_d = 12\text{ V}, R = 50\ \Omega, D = 0.8, r_L = 0.5\ \Omega$. Find ideal $V_o$, actual $V_o$, and efficiency $\eta$.

\tcbline
\textbf{Solution:}
\begin{enumerate}[leftmargin=*]
  \item Ideal: $V_{o,\text{ideal}} = \frac{V_d}{1 - D} = \frac{12}{0.2} = \mathbf{60\text{ V}}$.
  \item Parasitic factor: $\frac{r_L}{R (1 - D)^2} = \frac{0.5}{50 \times (0.2)^2} = \frac{0.5}{2.0} = 0.25$.
  \[
  V_o = \frac{V_d}{1 - D} \frac{1}{1 + 0.25} = 60 \times \frac{1}{1.25} = \mathbf{48\text{ V}}
  \]
  \item Efficiency: $\eta = \frac{1}{1 + 0.25} = \frac{1}{1.25} = \mathbf{80\%}$.
\end{enumerate}
\end{examplebox}

\begin{examplebox}[Problem 10: Buck-Boost Converter Sizing \& Device Stress]
\textbf{Question:} $V_d = 24\text{ V}, |V_o| = 36\text{ V}, R = 18\ \Omega, f_s = 50\text{ kHz}$. Find $D, I_L, L_{\min}$, and switch voltage stress.

\tcbline
\textbf{Solution:}
\begin{enumerate}[leftmargin=*]
  \item $\frac{|V_o|}{V_d} = \frac{D}{1 - D} \implies 1.5 = \frac{D}{1 - D} \implies D = \mathbf{0.60}$.
  \item Output current: $I_o = \frac{36}{18} = 2\text{ A} \implies I_L = \frac{I_o}{1 - D} = \frac{2}{0.4} = \mathbf{5\text{ A}}$.
  \item $L_{\min} = \frac{(1 - D)^2 R}{2 f_s} = \frac{(0.4)^2 \times 18}{2 \times 50,000} = \mathbf{28.8\ \mu\text{H}}$.
  \item Switch and Diode Peak Blocking Voltage:
  \[
  V_{\text{switch,peak}} = V_{\text{diode,PIV}} = V_d + |V_o| = 24 + 36 = \mathbf{60\text{ V}}
  \]
\end{enumerate}
\end{examplebox}

\newpage
%%%%%%%%%%%%%%%%%%%%%%%%%%%%%%%%%%%%%%%%%%%%%%%%%%%%%%%%%%%%%%%%%%%%%%%%%%%%%%%
\section{Master Formula Reference Summary}
%%%%%%%%%%%%%%%%%%%%%%%%%%%%%%%%%%%%%%%%%%%%%%%%%%%%%%%%%%%%%%%%%%%%%%%%%%%%%%%

\begin{table}[h!]
\centering
\small
\caption{Master Formula Sheet --- Phase Controlled Rectifiers}
\begin{tabularx}{\textwidth}{p{3.2cm}YYYY}
\toprule
\textbf{Converter Topology} & \textbf{Average Voltage $V_o$} & \textbf{RMS Voltage $V_{o,\text{rms}}$} & \textbf{Displ. Factor ($DF$)} & \textbf{Input Power Factor ($PF$)} \\
\midrule
1-$\phi$ Half-Wave ($R$) & $\frac{V_m}{2\pi}(1+\cos\alpha)$ & $\frac{V_m}{2}\sqrt{\frac{(\pi-\alpha)+\frac{1}{2}\sin 2\alpha}{\pi}}$ & --- & --- \\[1.5ex]
1-$\phi$ Full Converter & $\frac{2 V_m}{\pi}\cos\alpha$ & $V_s = \frac{V_m}{\sqrt{2}}$ & $\cos\alpha$ & $\frac{2\sqrt{2}}{\pi}\cos\alpha \approx 0.90\cos\alpha$ \\[1.5ex]
1-$\phi$ Semi-Converter & $\frac{V_m}{\pi}(1+\cos\alpha)$ & $\frac{V_m}{\sqrt{2}}\sqrt{\frac{(\pi-\alpha)+\frac{1}{2}\sin 2\alpha}{\pi}}$ & $\cos\left(\frac{\alpha}{2}\right)$ & $\frac{\sqrt{2}(1+\cos\alpha)}{\sqrt{\pi(\pi-\alpha)}}$ \\[1.5ex]
1-$\phi$ Full ($L_s$) & $\frac{2 V_m}{\pi}\cos\alpha - 2 f L_s I_o$ & --- & --- & $\Delta V_o = 2 f L_s I_o$ \\[1.5ex]
3-$\phi$ 6-Pulse Bridge & $\frac{3 V_{ml}}{\pi}\cos\alpha$ & $V_{ml}\sqrt{\frac{1}{2}+\frac{3\sqrt{3}}{4\pi}\cos 2\alpha}$ & $\cos\alpha$ & $\frac{3}{\pi}\cos\alpha \approx 0.955\cos\alpha$ \\[1.5ex]
3-$\phi$ Semi-Converter & $\frac{3 V_{ml}}{2\pi}(1+\cos\alpha)$ & --- & --- & Higher than 6-pulse \\[1.5ex]
3-$\phi$ Full ($L_s$) & $\frac{3 V_{ml}}{\pi}\cos\alpha - 3 f L_s I_o$ & --- & --- & $\Delta V_o = 3 f L_s I_o$ \\
\bottomrule
\end{tabularx}
\end{table}

\vspace{1em}

\begin{table}[h!]
\centering
\small
\caption{Master Formula Sheet --- DC-DC Switching Regulators (CCM)}
\begin{tabularx}{\textwidth}{p{3.2cm}YYYY}
\toprule
\textbf{Regulator Topology} & \textbf{Voltage Gain $\frac{V_o}{V_d}$} & \textbf{Inductor Current $I_L$} & \textbf{Boundary Inductance $L_{\min}$} & \textbf{Capacitor Ripple $\frac{\Delta V_o}{V_o}$} \\
\midrule
\textbf{Buck (Step-Down)} & $D$ & $I_o = \frac{V_o}{R}$ & $\frac{(1 - D) R}{2 f_s}$ & $\frac{1 - D}{8 L C f_s^2}$ \\[1.8ex]
\textbf{Boost (Step-Up)} & $\frac{1}{1 - D}$ & $\frac{I_o}{1 - D} = \frac{V_d}{(1 - D)^2 R}$ & $\frac{D (1 - D)^2 R}{2 f_s}$ & $\frac{D}{R C f_s}$ \\[1.8ex]
\textbf{Buck-Boost} & $-\frac{D}{1 - D}$ & $\frac{I_o}{1 - D} = \frac{D V_d}{(1 - D)^2 R}$ & $\frac{(1 - D)^2 R}{2 f_s}$ & $\frac{D}{R C f_s}$ \\
\bottomrule
\end{tabularx}
\end{table}

\vspace{1.5em}
\begin{center}
  \small\textcolor{gray}{--- End of Module 2 Master Study Guide ---}
\end{center}

\end{document}
